\begin{lemma}
\label{lemma-cover}
\begin{slogan}
Zariski-local properties of modules and algebras
\end{slogan}
Let $R$ be a ring. Let $M$ be an $R$-module. Let $S$ be an $R$-algebra.
Suppose that $f_1, \ldots, f_n$ is a finite list of
elements of $R$ such that $\bigcup D(f_i) = \Spec(R)$,
in other words $(f_1, \ldots, f_n) = R$.
\begin{enumerate}
\item If each $M_{f_i} = 0$ then $M = 0$.
\item If each $M_{f_i}$ is a finite $R_{f_i}$-module,
then $M$ is a finite $R$-module.
\item If each $M_{f_i}$ is a finitely presented $R_{f_i}$-module,
then $M$ is a finitely presented $R$-module.
\item Let $M \to N$ be a map of $R$-modules. If $M_{f_i} \to N_{f_i}$
is an isomorphism for each $i$ then $M \to N$ is an isomorphism.
\item Let $0 \to M'' \to M \to M' \to 0$ be a complex of $R$-modules.
If $0 \to M''_{f_i} \to M_{f_i} \to M'_{f_i} \to 0$ is exact for each $i$,
then $0 \to M'' \to M \to M' \to 0$ is exact.
\item If each $R_{f_i}$ is Noetherian, then $R$ is Noetherian.
\item If each $S_{f_i}$ is a finite type $R_{f_i}$-algebra, then
$S$ is a finite type $R$-algebra.
\item If each $S_{f_i}$ is of finite presentation over $R_{f_i}$, then
$S$ is a finitely presented $R$-algebra.
\end{enumerate}
\end{lemma}

\begin{proof}
We prove each of the parts in turn.
\begin{enumerate}
\item By Proposition \ref{proposition-localize-twice}
this implies $M_\mathfrak p = 0$ for all $\mathfrak p \in \Spec(R)$,
so we conclude by Lemma \ref{lemma-characterize-zero-local}.
\item For each $i$ take a finite generating set $X_i$ of $M_{f_i}$.
For each $x\in X_i$, write $x=y_x/f_i^{a_x}$ with
$y_x\in M$ and $a_x\geq0$, and put $Y_i=\{y_x\mid x\in X_i\}$.
Keep the original elements $x$: the image of $y_x$ is $f_i^{a_x}x$,
and $x=f_i^{-a_x}(y_x/1)$ in $M_{f_i}$.
Since $f_i$ is a unit in $R_{f_i}$, the images of $Y_i$ and the
original set $X_i$ generate the same $R_{f_i}$-submodule. Let $Y$
be the union of these sets. This is still a finite set.
Consider the obvious $R$-linear map $R^Y \rightarrow M$ sending the basis
element $e_y$ to $y$. By assumption this map is surjective after localizing
at an arbitrary prime ideal $\mathfrak p$ of $R$, so it is surjective by
Lemma \ref{lemma-characterize-zero-local}
and $M$ is finitely generated.
\item By (2) we have a short exact sequence
$$
0 \rightarrow K \rightarrow R^m \rightarrow M \rightarrow 0
$$
Since localization is an exact functor and $M_{f_i}$ is finitely
presented we see that $K_{f_i}$ is finitely generated for all
$1 \leq i \leq n$ by Lemma \ref{lemma-extension}.
By (2) this implies that $K$ is a finite $R$-module and therefore
$M$ is finitely presented.
\item By Proposition \ref{proposition-localize-twice}
the assumption implies that the induced morphism
on localizations at all prime ideals is an isomorphism, so we conclude
by Lemma \ref{lemma-characterize-zero-local}.
\item By Proposition \ref{proposition-localize-twice} the assumption
implies that the induced
sequence of localizations at all prime ideals is short exact, so we
conclude by Lemma \ref{lemma-characterize-zero-local}.
\item We will show that every ideal of $R$ has a finite generating set:
For this, let $I \subset R$ be an arbitrary ideal. By
Proposition \ref{proposition-localization-exact}
each $I_{f_i} \subset R_{f_i}$ is an ideal. These are all
finitely generated by assumption, so we conclude by (2).
\item For each $i$ take a finite generating set $X_i$ of $S_{f_i}$.
For each $x\in X_i$, write $x=y_x/f_i^{a_x}$ with
$y_x\in S$ and $a_x\geq0$, and put $Y_i=\{y_x\mid x\in X_i\}$.
The equalities $y_x/1=f_i^{a_x}x$ and
$x=f_i^{-a_x}(y_x/1)$ show that the images of $Y_i$ generate
the same $R_{f_i}$-subalgebra as $X_i$, because both scalar factors
belong to the base ring $R_{f_i}$. Let $Y$ be the union of the
finite sets $Y_i$; it is finite.
Consider the algebra homomorphism $R[X_y]_{y \in Y} \rightarrow S$
induced by $Y$. Since it is an algebra homomorphism, the image $T$
is an $R$-submodule of the $R$-module $S$, so we can consider the
quotient module $S/T$. By assumption, this is zero if we localize
at the $f_i$, so it is zero by (1) and therefore $S$ is an
$R$-algebra of finite type.
\item By the previous item, there exists a surjective $R$-algebra
homomorphism $R[X_1, \ldots, X_n] \rightarrow S$. Let $K$ be the kernel
of this map. This is an ideal in $R[X_1, \ldots, X_n]$, finitely generated
in each localization at $f_i$. Since the $f_i$ generate the unit ideal
in $R$, they also generate the unit ideal in $R[X_1, \ldots, X_n]$, so an
application of (2) finishes the proof.
\end{enumerate}
\end{proof}